% Paper size is configurable because the venue currently gives conflicting guidance.
\ifdefined\UseLetterPaper
\documentclass[conference,letterpaper]{IEEEtran}
\else
\documentclass[conference,a4paper]{IEEEtran}
\fi
\usepackage{amsmath,booktabs,graphicx,url}
\pagestyle{empty}
\title{Mapping Home Buyability: Payment Burden, Entry Capital, and Uncertainty Across U.S. Counties}
\author{
\IEEEauthorblockN{Ankit Hemant Lade}
\IEEEauthorblockA{Ascentt Systems Inc, USA\\
ankitlade12@gmail.com}
\and
\IEEEauthorblockN{Venkata Rao Kondepati}
\IEEEauthorblockA{Ascentt, USA\\
venkata.kondepati@gmail.com}
}
\begin{document}
\maketitle
\thispagestyle{empty}
\begin{abstract}
A county can appear affordable on a mortgage-payment map while requiring substantial cash at purchase. We examine how measurement choices change county comparisons by combining Zillow home-value estimates, American Community Survey income and ownership-cost data, and explicit financing scenarios. A reproducible Python and DuckDB pipeline produces 18,372 county-scenario records. An interactive explorer links county maps with a payment-versus-capital view and distinguishes missing observations, income uncertainty, and unbounded cost estimates. On a common sample of 3,036 counties, replacing all-household income with renter income moves 1,682 counties above a 28\% payment benchmark. Adding lower-bound tax and insurance estimates moves another 424 above it. Only 65.8\% of counties in the original highest-burden decile remain there after both changes. Lower down payments reduce upfront capital requirements while increasing recurring costs. A separate 2019--2024 accounting analysis describes changes associated with nominal prices, incomes, and mortgage rates. These findings show how the choice of income population and cost measure alters the geography of apparent affordability. The analysis evaluates representations of county-level conditions; it does not establish individual mortgage eligibility or causal policy effects.
\end{abstract}
\begin{IEEEkeywords}
Housing affordability, visual analytics, geospatial analysis, data integration, uncertainty visualization.
\end{IEEEkeywords}

\section{Introduction}
Buying a home requires both a manageable monthly payment and enough cash to complete the purchase. These requirements can pull in opposite directions: a smaller down payment lowers the cash needed at closing but increases the loan balance and may add mortgage insurance. A map based only on principal and interest captures one side of this tradeoff.

We use the term home buyability to describe the capacity to enter ownership. County data capture only part of this capacity. They nevertheless allow us to compare how different measures describe access to ownership across places. Two decisions are particularly consequential. The first is whether to compare housing costs with the income of all households or renter households. The second is whether to include ownership costs beyond principal and interest. Income uncertainty and incomplete observations add a further complication near any affordability threshold.

This study asks how these choices change county rankings and benchmark classifications, and how a visual interface can expose those differences without hiding gaps in the data. We combine a reproducible data pipeline with paired maps, a payment-versus-capital view, and explicit missing-value and uncertainty states. The evaluation measures changes in the resulting comparisons. It does not test whether the interface improves users' decisions.

\section{Related work}
The National Association of Realtors Housing Affordability Index already combines prices, income, and mortgage financing under explicit assumptions \cite{nar}. We use this established relationship as a baseline. HomeSeeker integrates heterogeneous location-centered real-estate data and supports interactive property and suburb exploration \cite{homeseeker}. Its focus on property exploration provides relevant context for our county-level comparison of financing and income assumptions. Lucchesi and Wikle demonstrate several ways to visualize uncertainty in areal estimates using ACS data and margins of error \cite{lucchesi}. Their work motivates treating uncertainty as part of the mapped information.

We build on these approaches by comparing recurring costs and entry capital under a common set of county data and financing assumptions. The contribution is the reproducible comparison and its empirical results, rather than a new mortgage formula or visual encoding.

% Declare the wide map early so IEEE two-column layout places it on the next page.
\begin{figure*}[t]
\centering
\includegraphics[width=\textwidth]{../outputs/payment_burden_pilot.pdf}
\caption{County principal-and-interest burdens using all-household and renter income under identical financing terms. The scale is fixed across maps; gray indicates missing data. Alaska and Hawaii are included in the analysis but omitted from these maps.}
\label{fig:county-maps}
\end{figure*}

\section{Data and measures}
The geographic universe comprises 3,144 county equivalents in the 50 states and District of Columbia using 2024 Census cartographic boundaries. We average twelve monthly 2024 Zillow Home Value Index observations \cite{zhvi} and join ACS 2020--2024 five-year estimates \cite{acs} of median household income and renter-household income. ZHVI is a typical-value measure, not the median asking price of homes available for sale. There are 3,062 valid price/all-household-income matches and 3,043 renter-income matches. Missing Zillow coverage and incompatible Connecticut county/planning-region definitions are retained in the coverage audit rather than silently imputed.

For home value $V$, down-payment share $d$, annual mortgage rate $i$, and annual income $Y$, monthly principal and interest and its income ratio are
\begin{equation}
\begin{aligned} M&=V(1-d)\frac{r}{1-(1+r)^{-360}},\\ r&=i/12,\qquad B_{PI}=12M/Y.\end{aligned}
\end{equation}
The baseline uses $d=0.20$ and $i=0.0703$, the September 24, 2026 PMMS observation \cite{pmms}. This is a fixed financing scenario applied to 2024 inputs, not an estimate of contemporaneous 2026 county conditions. We compare 5\%, 10\%, and 20\% down payments and a separately identified 0.50-percentage-point rate stress. These yield 18,372 county-scenario records.

ACS B25103 supplies annual tax medians for mortgaged owner-occupied homes. B25141 supplies insurance-cost bins \cite{insurance}. We retain tax censoring and derive an interval containing the insurance median from the central observations' bins. Open upper categories remain unbounded. Recurring cost proxies add these components to P\&I and, below 20\% down, a PMI assumption of \$30--\$70 monthly per \$100,000 of principal. Upfront cash assumes the down payment plus 2--5\% closing costs, following Freddie Mac's consumer ranges \cite{pmi,closing}. These are scenario ranges, not borrower-specific quotes or confidence intervals. Adding separate county medians does not recover the observed median total cost.

We transform income plus/minus its reported ACS 90\% margin of error \cite{acsmoe} while holding price fixed to classify whether the P\&I ratio is below, above, or straddles 28\%. This reference supports consistent comparisons; it does not determine mortgage approval. We do not observe assets, other debts, credit scores, assistance, maintenance, HOA charges, or property-specific reassessment and insurance repricing.

% Declare the results figure before its discussion to avoid a final-page float.
\begin{figure*}[t]
\centering
\includegraphics[width=\linewidth]{../outputs/representation_comparison.pdf}
\caption{Paired comparisons on a common county sample. Fixed plot bounds retain outlying counties in the statistics; the figure reports the clipping count.}
\label{fig:representation-comparison}
\end{figure*}

\section{Implementation and evaluation}
We use Python and DuckDB to join the source tables and calculate the scenarios. Each cached input has a recorded URL, retrieval metadata, and SHA256 hash. An audit verifies all 163 source-file hashes against a frozen manifest, checks unique scenario keys, reconciles SQL and extended payment calculations, verifies historical additivity, and preserves nullable censoring flags. Numerical agreement across implementations demonstrates internal consistency, not external validity. An offline HTML explorer provides coordinated county maps, a cash-versus-payment scatterplot, scenario controls, county inspection, and CSV export. Fixed map bins support comparisons. Figure~\ref{fig:county-maps} shows the difference between the two income populations under the baseline financing scenario. Missing values and unbounded upper costs receive distinct encodings. Display maps cover the contiguous United States; county selection and statistics include Alaska and Hawaii. Simplified display geometry is not used for analytical boundary checks.

The workflow also exports county results as GeoJSON for inspection in external GIS software. The evaluation reported here concerns the standalone explorer.

We compare ranks using Spearman correlation, highest-decile set overlap, and threshold transitions on common samples. Decile size is rounded upward, and GEOID breaks boundary ties. Threshold sensitivity includes 25\%, 28\%, 30\%, and 36\%. These evaluate consequences of representation choices, not predictive accuracy against individual ownership outcomes.

\section{Results}
Figure~\ref{fig:representation-comparison} compares the original payment ratio with the renter-income and ownership-cost measures.
\begin{table*}[t]
\caption{Representation comparisons at 7.03\% and 20\% down. Sample differences reflect missing cost inputs.}
\centering\small
\begin{tabular}{lrrr}\toprule
Comparison & Counties & Rank correlation & Top-decile retained\\\midrule
All-household to renter income & 3,043 & 0.868 & 71.1\%\\
Renter P\&I to lower cost proxy & 3,036 & 0.981 & 88.8\%\\
Original to renter lower cost proxy & 3,036 & 0.818 & 65.8\%\\\bottomrule
\end{tabular}
\end{table*}
On the common sample of 3,036 counties, changing from all-household to renter income moves 1,682 counties from at or below 28\% to above it. Adding the lower-bound ownership-cost estimates moves another 424 counties above the benchmark, bringing the combined count to 2,106. Adding positive costs must raise the ratio; the empirical question is how much this changes classifications and county rankings. Using the larger income-only sample of 3,043 counties gives 1,685 upward crossings, which explains the difference between that comparison and the common-sample result.

For 3,036 common renter-income/cost counties, 20\% down produces median county recurring-cost ratios of 44.5--45.7\% and upfront cash ratios of 126.9--144.2\% of annual renter income. At 5\% down, the corresponding values are 53.6--57.3\% and 40.4--57.7\%. Lower down payments reduce the initial capital requirement while increasing recurring burdens. The cash ratios are not years required to save, because saving rates and assets are unobserved.

For the original renter P\&I scenario, income-MOE sensitivity leaves 1,924 counties above 28\% throughout, 258 below, and 861 straddling. Nineteen baseline counties lack valid inputs for this calculation; an additional 82 geographic-universe counties lack the baseline price/income match. Thus, the full explorer reports 101 missing classifications. These are income-only ranges, not simultaneous confidence statements across counties.

\subsection{Historical accounting}
We compare annual-average 2019 and 2024 ZHVI and PMMS rates with the corresponding ACS five-year income releases. A geographic screen retains counties with at most 1\% relative symmetric-difference area between boundary vintages. This is a boundary-change exclusion rule, not full geographic harmonization. The matched historical samples contain 2,925 counties with all-household income data and 2,912 with renter-income data. The historical renter sample contains 90.74\% of 2024 renter households in the geographic universe, compared with 98.85\% for the 3,036-county representation sample. These are sample-coverage measures based on ACS household counts, not percentages of households able to buy.

Median all-household P\&I burden rises from 13.91\% to 21.42\%; renter burden rises from 22.94\% to 35.91\%. A six-order Shapley accounting decomposition allocates the mean all-household increase of 8.23 percentage points into nominal prices (+7.14), nominal income ($-4.83$), and rates (+5.92). Renter contributions are +11.67, $-7.83$, and +9.72 points, summing to +13.55. Exact additivity is checked. Allocations use endpoint nominal dollars and are neither causal effects nor isolated inflation effects. PMMS methodology changed between endpoints \cite{pmmsmethod}.

\section{Limitations}
A ratio of county medians is not a distribution of household eligibility. Renter households are a relevant comparison population, but not a sample of actual prospective purchasers; selecting them does not validate a first-time-buyer interpretation. Property values and household income are not jointly observed for a matched purchaser/property pair. Counties aggregate heterogeneous neighborhoods, and conclusions may change with geographic scale (the modifiable areal unit problem). Existing-owner taxes and insurance may differ substantially from new-purchaser costs. The ZHVI series is model-derived and retrospectively revised; historical values here come from the downloaded vintage, not the information set available in 2019. ACS estimates pool five years; aligning their release year with annual home values does not create synchronous observations. Excluded geographies and boundary screens affect generalizability. A common national rate change mechanically rescales P\&I ratios and preserves rankings with other inputs fixed; its map cannot identify heterogeneous causal rate effects.

The analysis does not estimate rental displacement, the effect of mortgage lock-in on housing inventory, or a separate effect of inflation. The reported ranges account for income uncertainty and cost categories. They do not incorporate sampling error in the tax and insurance estimates or uncertainty in ZHVI. Insurance bounds describe bins of estimated counts, not the uncertainty of the population median. Rank and threshold summaries are descriptive point-estimate comparisons; statistical significance of changes is not claimed. No participant study establishes usability or improved comprehension. Statistical metropolitan-core/periphery comparisons remain outside the completed analysis.

\section{Conclusion}
The same counties can look substantially different when housing costs are compared with renter income and when taxes and insurance are included. Upfront cash requirements reveal a second constraint that a monthly-payment map cannot show. Presenting these measures together, while retaining missing values and uncertainty, makes the assumptions behind each comparison visible. The resulting analysis describes county-level exposure under specified financing terms; assessing which households can actually buy requires household resources and available-property data that this study does not observe.
\bibliographystyle{IEEEtran}
\bibliography{references}
\end{document}
